\section{Part C. Bounding and Optimization}
\begin{frame}{Branch-and-Bound의 정의}
\ssmission{최적화 state-space tree에서 optimistic objective bound를 계산하고, incumbent를 개선할 수 없는 subtree를 prune한다.}
\begin{description}
\item[Incumbent] 지금까지 찾은 best feasible solution.
\item[Live node] frontier에 있어 향후 확장할 node.
\item[Expanded node] frontier에서 꺼내 child를 만든 node.
\item[Dead node] infeasible, bounded, 또는 확장이 끝난 node.
\end{description}
\end{frame}
\begin{frame}{Backtracking vs Branch-and-Bound}
\small
\begin{columns}[T]
\column{.48\linewidth}
\begin{block}{Backtracking}
\begin{itemize}
\item 목표: feasible solution 탐색
\item prune 증거: feasible descendant 없음
\item 핵심: feasibility / \textbf{promising}
\item 흔히 DFS지만 correctness는 DFS에 의존하지 않음
\end{itemize}
\end{block}
\column{.48\linewidth}
\begin{block}{Branch-and-Bound}
\begin{itemize}
\item 목표: optimal solution 탐색
\item prune 증거: incumbent 개선 불가
\item 핵심: objective, incumbent, safe bound
\item FIFO, LIFO, best-first 모두 가능
\end{itemize}
\end{block}
\end{columns}
\begin{alertblock}{용어 정책}이 강의에서 \textbf{promising}은 feasibility 가능성만 뜻한다. 최적화 prune은 ``bound can improve incumbent''로 표현한다.\end{alertblock}
\end{frame}
\begin{frame}{Max와 Min의 Prune 부등호}
\begin{columns}
\column{.48\linewidth}
\begin{block}{Maximization}
\[
UB(node)\ge OPT_{\mathrm{desc}}(node)
\]
\[
UB(node)\le INC\Rightarrow P
\]
\end{block}
\column{.48\linewidth}
\begin{block}{Minimization}
\[
LB(node)\le OPT_{\mathrm{desc}}(node)
\]
\[
LB(node)\ge INC\Rightarrow P
\]
\end{block}
\end{columns}
\begin{exampleblock}{Safe-bound proof}Prune 조건이 성립하면 모든 feasible descendant도 incumbent보다 좋을 수 없으므로, 더 나은 해는 제거되지 않는다.\end{exampleblock}
\end{frame}
\begin{frame}{Optimistic Bound가 필요한 이유}
\begin{itemize}
\item Maximization UB가 실제 가능한 최고값보다 작으면 optimal branch를 잘못 제거할 수 있다.
\item Minimization LB가 실제 가능한 최저값보다 크면 같은 문제가 생긴다.
\item loose하지만 safe한 bound는 correct, tight하지만 unsafe한 bound는 incorrect.
\end{itemize}
\end{frame}
\begin{frame}{Maximization Best-First Skeleton}
\tiny\begin{algorithmic}[1]
\Procedure{BranchAndBound}{$root$}
\State \(incumbentValue\gets-\infty;\ PQ\gets\) max-heap by upperBound
\State compute \(root.upperBound\); insert \(root\)
\While{\(PQ\ne\varnothing\)}
\State \(node\gets ExtractMax(PQ)\)
\If{\(node.upperBound\le incumbentValue\)}\State\textbf{continue}\EndIf
\If{\Call{RepresentsFeasibleSolution}{node} and \(node.value>incumbentValue\)}
\State update incumbent
\EndIf
\If{\Call{IsComplete}{node}}\State\textbf{continue}\EndIf
\ForAll{\(child\in Expand(node)\)}
\If{\Call{FeasibilityPossible}{child}}\State compute bound
\If{\(child.upperBound>incumbentValue\)}\State insert child\EndIf\EndIf
\EndFor
\EndWhile
\State\Return incumbent
\EndProcedure
\end{algorithmic}
\end{frame}
\begin{frame}{Bound-Based Pruning Animation}
\centering\begin{tikzpicture}
\node[sscurrent](r)at(0,2){UB 20};\node[ssopen](a)at(-2,.5){UB 16};\node[ssopen](b)at(2,.5){UB 12};
\draw[ssedge](r)--(a);\draw[ssedge](r)--(b);
\only<2-|handout:1>{\node[ssvalid](f)at(-2,-1){value 14};\draw[ssedge](a)--(f);}
\only<3-|handout:1>{\node[ssinc]at(0,-2){INC = 14};}
\only<4-|handout:1>{\node[ssbound]at(2,.5){P};\node[font=\scriptsize,text=violet]at(3.5,.5){UB 12 \(\le14\)};}
\node[ssnote]at(0,-3){\only<1|handout:0>{root 확장}\only<2|handout:0>{feasible solution 발견}\only<3|handout:0>{incumbent 갱신}\only<4|handout:1>{더 좋아질 수 없는 B subtree prune}};
\end{tikzpicture}
\end{frame}
\begin{frame}{Bound Quality}
\[
\text{total work}\approx
\text{expanded nodes}\times\text{bound cost}.
\]
\small\begin{tabular}{lll}\toprule
Bound&장점&단점\\\midrule
loose&계산이 빠름&pruning 적음\\
tight&pruning 많음&계산이 비쌀 수 있음\\\bottomrule
\end{tabular}
\end{frame}
\begin{frame}{Branch-and-Bound Takeaway}
\sstakeaway{Branch-and-Bound는 탐색 순서가 아니라 incumbent와 safe objective bound의 관계로 정의된다.}
\end{frame}
